\documentclass[11pt]{article}
\usepackage[margin=0.8in]{geometry}
\usepackage[colorlinks=true,urlcolor=blue]{hyperref}
\usepackage{microtype}
\pagestyle{empty}
\setlength{\parindent}{0pt}
\setlength{\parskip}{0.7em}

\begin{document}

{\Large \textbf{Yitian Wang}}\\
+1(929)624-1021 \quad
\href{mailto:ywang00149@gmail.com}{ywang00149@gmail.com} \quad
\href{https://cosmotim.github.io}{cosmotim.github.io}\\
\href{https://linkedin.com/in/tim-wang-yitian}{linkedin.com/in/tim-wang-yitian}

\vspace{1em}
September 15, 2026

Hiring Committee\\
Exponent, Thermal Sciences Practice\\
Natick, MA

\textbf{Re: Thermal Sciences Associate (Ph.D.)}

Dear Hiring Committee,

I am applying for the Thermal Sciences Associate position at Exponent. My Ph.D. is in Electrical and Computer Engineering, but my research centers on the thermal transport, thermodynamics, and electrochemistry of solid-state energy materials. I am drawn to Exponent's work because it applies rigorous physical analysis and experimental evidence to high-consequence engineering questions.

At UC Riverside, I designed and executed experimental programs to understand heat flow in lithium-ion solid electrolytes. I developed a floating-zone crystal-growth workflow that produced centimeter-scale LLZTO single crystals and enabled phonon-resolved neutron-scattering measurements. I then combined these measurements with thermal-transport characterization and a two-channel phonon--diffuson model in MATLAB to explain anomalous temperature-dependent behavior that conventional models could not capture.

As experiment lead at Oak Ridge National Laboratory, I authored two successful beam-time proposals and led inelastic neutron-scattering experiments at HFIR and SNS. I made real-time experimental decisions, built automated Mantid and Python data-reduction pipelines, and coordinated work across approximately seven research groups. Earlier, at Columbia University, I fabricated composite battery separators, assembled and tested CR2032 cells, and used EIS and galvanostatic cycling to evaluate electrochemical performance and interfacial stability.

These experiences have prepared me to plan experiments, work directly with instrumentation and complex data, develop practical models, and communicate technical conclusions clearly. My publication record includes six first-author and three co-authored peer-reviewed journal articles, including a recent review in \textit{Applied Physics Reviews}. I would welcome the opportunity to bring this combination of thermal-sciences research, experimental judgment, and technical writing to Exponent's client-facing consulting work.

Thank you for your consideration. I look forward to discussing how my background can contribute to the Thermal Sciences Practice.

Sincerely,\\[1.5em]
Yitian Wang

\end{document}
